\documentclass[11pt]{article}
\usepackage{enrico}
\usepackage{booktabs}
\usepackage{enumitem}
\usepackage[hidelinks]{hyperref}
\setlist{itemsep=2pt,topsep=3pt}

\title{Mill Street Labs, Sandbox 1: the research program\\[2pt]
\large What we test, why it matters for the company, and how we keep it rigorous}
\author{Draft prepared with Claude for Enrico Yao-Bate}
\date{October 6, 2026 (status: nothing below has been run on a real model; see \S\ref{sec:status})}

\begin{document}
\maketitle

\section{What this document is}

A single place that states (i) the company's goal and edge, (ii) the game we use and every term we use about it,
(iii) the statistical framework the experiments live in, (iv) the experiments themselves, each written as
\emph{what it is for, what it is, why it is worth doing, how it is done}, (v) the rigor protocol every experiment
follows, and (vi) the build status and the next steps. It is meant to be read by someone who knows statistics but
not Hanabi, reinforcement learning or causal inference, so terms are defined before use. Claims about prior work are
phrased as ``not found in our search of date X'' with the nearest works named; see \S\ref{sec:novelty}.

\section{Goal, hypothesis, edge}\label{sec:goal}

\paragraph{Goal.} Build the layer that lets a population of agents discover knowledge, share it and keep building on
it over long horizons. The product is the organization: who teaches whom, how a learner checks what it is taught, how
credit and compute flow. Sandbox~1 (this document) is the controlled laboratory; Sandbox~2 (market making) points the
same layer at the real domain.

\paragraph{Hypothesis.} Collective improvement comes from \emph{adaptive social dynamics} (groups compete for finite
resources so helping the group pays; agents revise how they cooperate and migrate) and \emph{intelligence transfer}
(teachers are credited for students' improvement; discoveries from different lineages are combined; learners verify
before adopting).

\paragraph{Edge, as stated in the literature map.} Five things not found together in prior work: (1) the group as
the unit that competes and reproduces, with teaching; (2) causal credit over a teaching DAG with several teachers;
(3) controlled perturbation of organization over long runs; (4) a formal model predicting when organization beats
scale; (5) designing against persuasive-but-wrong ideas and against convergence. Rigor itself is part of the edge:
the adjacent literature runs single seeds, no counterfactuals, no pre-registration.

\paragraph{What preseed has to show.} Transfer alone is not interesting; many systems transfer knowledge. The cheap
edge items are (2), (4) and (5), and the deliverable is one or two defensible figures for them plus the methods that
produced them. Items (1) and (3) need a factorial over organization rules and wait for funding.

\section{The game, from scratch}\label{sec:game}

Hanabi is a \emph{cooperative} card game. We play it with two players who are on one team against the deck; nobody in
the game competes with anybody. Groups and competition exist only in how we organize agents outside the game.

\paragraph{Goal.} Build five fireworks, one per colour, by playing cards onto each colour's stack in the order
$1,2,3,4,5$. Score is the number of cards successfully played, maximum 25.

\paragraph{The twist.} You hold your cards facing away from you: you see your partner's hand, never your own. The
only way to learn what you hold is from clues your partner gives you.

\paragraph{Components.} 50 cards: five colours, ranks 1 to 5, with three 1s, two each of 2, 3, 4 and one 5 per
colour. Five cards in hand each. Eight clue tokens and three life tokens, shared.

\paragraph{A turn is one of three actions.} \emph{Clue:} spend a token and tell your partner one colour or one rank,
pointing at every card in their hand that matches. \emph{Discard:} put one of your own cards face up, regain a token,
draw. \emph{Play:} put one of your own cards on the stacks; if it is the next rank of its colour it stays and scores,
otherwise it is a bomb and costs a life; draw either way.

\paragraph{Ending.} All stacks complete (25); or three bombs (our engine scores this 0); or the deck runs out and
each player takes one last turn.

\paragraph{Why conventions exist.} Eight tokens cannot describe 25 cards literally, so good teams agree in advance on
rules that make clues mean more than they say. Example: the partner's newest card is a red~1; you clue ``red'',
touching only that card. Literally: ``this card is red''. Under the common convention ``a clue touching one unknown
card means play it'', the partner plays it without knowing its rank. If the partner instead believes ``a red clue means
save it'', they hold it forever. Same clue, two outcomes, depending only on whether the rule is shared. The human
community on \texttt{hanab.live} maintains a written rulebook of conventions in levels 1 to 25; experts score 23 to 24.

\paragraph{Why this game for us.} Strategies are explicit rules, writeable as text and code. A rule has value only if
the partner shares it, so knowledge must be transmitted to be worth anything. Compatibility is measurable by playing
two bots together. The score is absolute with human and RL reference points. A game simulates in a millisecond.

\section{Objects and notation}\label{sec:objects}

\begin{itemize}
\item \textbf{Deal.} One shuffled deck order, indexed by a seed. \textbf{Paired seeds:} every bot in a comparison
plays identical deals, so differences are due to the bots, not the draw.
\item \textbf{Policy, or bot,} $\pi$. A deterministic program mapping what one player can see to a legal move.
\item \textbf{Score functional.} $V(\pi,\pi')$ is the expected score when $\pi$ plays with $\pi'$ over the deal
distribution. $V(\pi,\pi)$ is \emph{self-play}; $V(\pi,\pi')$ for $\pi\ne\pi'$ is \emph{cross-play}, the
compatibility measure. Known phenomenon (Hu et al.\ 2020): independently trained bots can score $24$ with
themselves and $2.5$ with each other.
\item \textbf{Anchors.} Fixed reference bots with known scores (random: 0; Canaan et al.'s Piers: 16.99; Flawed: 0),
so results are comparable across experiments and time.
\item \textbf{Conventions document.} Plain-English rules the bot follows. \textbf{Artifact} $a=(\text{code},
\text{conventions})$, with an id. The unit of knowledge.
\item \textbf{Agent.} An LLM-driven process owning one current artifact, its \emph{incumbent}, able to revise it,
send it, and receive others'.
\item \textbf{Message} $m=(a,e)$ from teacher $i$ to student $j$: an artifact plus \emph{evidence} $e$, the scores
our harness measured for $a$ (signed by the harness; agents cannot fabricate it).
\item \textbf{Verification.} Before adopting, $j$ plays $N$ deals with $a$ and checks it beats its incumbent.
\textbf{Adoption:} $j$ switches incumbent.
\item \textbf{Corpus} $\C_t$: the store of artifacts with evidence; \textbf{touch log:} who retrieved what, when.
\item \textbf{Generation.} One round in which every agent receives, verifies, revises, evaluates and teaches.
\item \textbf{Population state} $(P_t,\C_t)$: incumbents and corpus. \textbf{Group:} a subset with more exchange
inside than between.
\item \textbf{Student update.} $\pi_j^{t+1}=U(\pi_j^t,\,W,\,\xi)$ where $W\in\{0,1\}^k$ says which of $k$ messages
were delivered, $U$ contains verification, adoption and the LLM revision, and $\xi$ is the LLM's sampling
randomness. On a local model with a fixed sampling seed, $\xi$ is fixed and $U$ is a deterministic function.
\end{itemize}

\section{The statistical framework}\label{sec:framework}

\subsection{Potential outcomes for one student}

Following Wager (\emph{Causal Inference: A Statistical Learning Approach}, ch.~1, pp.~3--16): the unit is a
(student, generation); the treatment is the delivered set $W$; the outcome $Y(W)$ is the student's held-out
self-play improvement, $V(\pi',\pi')-V(\pi,\pi)$ on the same deals (paired). The average main effect of message $j$
is $\tau_j=\E_W[Y(W\cup\{j\})-Y(W\setminus\{j\})]$. SUTVA (Wager eq.~1.5, p.~5) holds along one student's
trajectory by construction because $U$ is deterministic; and the delivery probabilities are \emph{known}, since we
set them, which is the favourable case of Wager's Corollary~3.3 (p.~38): with the true propensity known, the
augmented IPW estimator is efficient with any consistent regression model. Regression adjustment (here, blocking on
the deal seed) is never worse than the raw difference in means (Wager Thm~1.3, eq.~1.20, p.~14), so we always pair.

\subsection{Credit as a game-theoretic quantity}

With $k$ messages and a value function $v(S)=Y(S)$ for delivered subset $S$, the Shapley value
$\phi_j=\sum_{S\not\ni j}\frac{|S|!\,(k-|S|-1)!}{k!}\,[v(S\cup\{j\})-v(S)]$ is the credit formalism of Yang et
al.\ 2025 (arXiv 2511.10687, ``one message is one player''), which we adopt and implement ourselves (no code was
released). It satisfies efficiency, $\sum_j\phi_j=v(\text{all})-v(\emptyset)$; the main effect $\tau_j$ does not.
Because $U$ is deterministic and seeded, exact replay over all $2^k$ subsets costs $2^k$ LLM calls and gives
$\phi$ exactly; at $k\le 7$ this is affordable as an oracle. Cheap estimators: leave-one-out ablation as in SHARP
(arXiv 2602.08335), $k+1$ replays, first order only; singles-plus-pairs replay as in Li et al.\ 2026 (arXiv
2608.29228), $1+k+\binom{k}{2}$ replays, recovers pairwise interactions; and randomized delivery plus regression
(a datamodels-style estimator), zero extra replays, amortized across the population.

\subsection{Interference between agents}

Within a group, one agent's treatment changes what is available to others, so SUTVA fails across agents. Wager
ch.~11--12 (pp.~144--171) gives the tools: an exposure mapping $H_i(w)$ (Assumption~11.1, p.~146); the network
interference model where $i$'s outcome depends on treatments in $\{i\}\cup N_i$ (Def.~11.1, p.~147); the IPW exposure
estimator, unbiased under Bernoulli randomization with known exposure probabilities (Thm~12.1, p.~157); the
finite-population Neyman variance, which needs no superpopulation of agent cohorts (Thm~12.2, p.~159); and the
HAC variance built on the randomization dependency graph $G_{ij}$ (eq.~12.18, p.~163), made conservative by
projecting $G$ to be PSD (Cor.~12.6, p.~167). The permutation test for the sharp null (Thm~11.1, p.~149) and for
``no spillovers'' (Thm~11.2, p.~151--152) are what we use to \emph{test} for interference before assuming it away.
Inference stays valid while any unit's reach grows slower than $n^{1/4}$ (Thm~12.7, p.~168), which caps how many
students one teacher may touch per generation relative to population size.

\subsection{Evaluating an organization rule from logs}

A routing or verification rule is a policy over deliveries across generations. Evaluating a rule from logs
collected under another rule is off-policy evaluation. Sutton and Barto (ch.~5, pp.~103--115): the coverage
assumption (p.~103), the importance ratio (eq.~5.3, p.~104), ordinary versus weighted importance sampling
(eqs.~5.5--5.6, p.~104), with weighted ``strongly preferred'' because its variance is bounded (p.~105), and
per-decision weighting to drop reward-irrelevant tail factors (eqs.~5.9--5.10, 5.15). Over many generations the
sequential weights compound (Wager eq.~14.12--14.14, pp.~195--196, the curse of horizon); if the student state
plausibly mixes, model it as an MDP and use the doubly robust stationary estimator (Wager eq.~15.9, p.~208; CLT
Thm~15.3, p.~211). If a rule is itself chosen adaptively over generations, sample means are biased (Wager
pp.~76--77); use the adaptively weighted estimator (eq.~6.12, Thm~6.2, p.~78) or a switchback design with bias
$\le 4M\lambda(1+t_0)$ (Thm~15.5, p.~223).

\subsection{Credit over time and teacher choice}

Minsky's credit-assignment problem (``how do you distribute credit for success among the many decisions that may
have been involved in producing it'', quoted in Sutton and Barto p.~17) is our problem with ``decisions'' replaced
by ``messages''. Crediting a teacher for a student's improvement several generations later is the backward view of
TD($\lambda$): an eligibility trace per (teacher, student) decaying by $\gamma\lambda$ per generation (eqs.~12.5--12.7,
p.~292--294), with $\lambda=1$ recovering Monte Carlo credit when bootstrapping is not trusted (p.~294). Choosing
which teacher to listen to is a nonstationary bandit: constant step size $Q_{n+1}=Q_n+\alpha(R_n-Q_n)$ (eq.~2.5,
p.~32) rather than sample averages, with the book's own caveat that UCB needs more than that under nonstationarity
(p.~36). Any baseline that does not depend on the action leaves the expected update unchanged but reduces its
variance (p.~329); and shaping rewards can change behaviour in unintended ways (p.~470), which is why organization
rewards here are measured outcomes, not bonuses.

\section{Rigor protocol}\label{sec:rigor}

Every experiment follows this before it runs.
\begin{enumerate}
\item \textbf{Pre-registration} (one page): question; setting; estimand; planted or known truth where possible;
design (what is paired, what is randomized, what is blocked); metrics; decision rules with thresholds fixed now;
minimum effect worth detecting; a pilot of 5 seeds to estimate variance, then the seed count by a power calculation;
stopping rule; what would falsify the approach; compute and cost. Stored in \texttt{files/prereg/} and logged in
the ledger.
\item \textbf{Paired seeds} for every comparison; three disjoint deal sets per generation (feedback traces shown to
the LLM; verification; held-out evaluation), with the generalization gap reported.
\item \textbf{Null populations.} The whole pipeline run with a null stub that only adds noise, giving the null
distribution of every reported metric; every figure carries a null band.
\item \textbf{Evidence ladder.} Stub, then local model A, then local model B as a blocking factor; a result is
accepted only if the pre-registered effect's interval excludes zero at the declared minimum effect in both blocks.
\item \textbf{Multiplicity.} Benjamini--Hochberg at $q=0.1$ across all confirmatory tests in the ledger; negatives
stay in the ledger.
\item \textbf{Novelty and reuse.} Before a claim is written, a targeted search and the three nearest works
(\S\ref{sec:novelty}). Existing tools are adopted when they fit; novelty only where nothing fits.
\item \textbf{Cost accounting.} Every curve is plotted against tokens and dollars by call type; local runs are
\$0 but seconds are logged.
\end{enumerate}

\section{The experiments}\label{sec:experiments}

Each experiment: \emph{For} (which goal or edge item), \emph{What} (plain description), \emph{Why} (why it is worth
doing now), \emph{How} (design and estimator), \emph{Decision rule}, \emph{Cost}, \emph{Status}. All of C1, D1,
A1, B1 share one setup, one student with a few teachers and many replicate seeds, so the harness work is shared.

\subsection{C1. Which credit estimator to trust (edge item 2)}\label{sec:c1}
\textbf{For.} The incentive mechanism of the one-pager: teachers credited for students' improvement.\\
\textbf{What.} A student receives three messages whose true effects we know because they carry anchor artifacts
(Piers, large positive; Flawed with persuasive prose, negative or zero depending on verification; the student's own
incumbent, zero). We compute the exact Shapley credit by replaying all eight delivery subsets, then ask which cheap
estimator recovers it.\\
\textbf{Why.} Credit is what makes teaching in an agent's interest, and the adjacent literature either has no
empirical validation (Yang et al.) or no code (all four nearest papers). An estimator we can trust every generation
is a prerequisite for everything organizational.\\
\textbf{How.} Estimand $\tau_j$ and $\phi_j$ as in \S\ref{sec:framework}. Estimators: leave-one-out ablation
($k+1$ replays), equal-split paired delta (1 replay), randomized delivery plus ridge regression (0 extra replays),
and in a second arm with $k=7$ (four null teachers added) singles-plus-pairs replay versus an 8-run Plackett--Burman
design versus 8 Bernoulli runs, scored against an exhaustive 128-run oracle on 5 seeds. Two strata: verification on
and off (intent-to-treat: the adoption rule is inside the effect). Pilot 5 seeds, then $R$ by power, capped at 60;
replication block on a second local model.\\
\textbf{Decision rule.} Accept an estimator if Spearman rank correlation with the oracle has IQM $\ge0.9$ (bootstrap
lower bound $\ge0.8$), RMSE $\le2$ points, sign error on the sabotaged teacher $\le5\%$, in both blocks; prefer the
cheapest accepted. Primary contrast: RMSE(randomized delivery) minus RMSE(leave-one-out), minimum effect 1 point.
If nothing passes at $R=60$: causal credit is not available at this cost and the one-pager's claim is reworded.\\
\textbf{Cost.} $8\times R\times 2$ local calls; at $R=30$ about 12 hours on the M4 Pro; \$0.\\
\textbf{Status.} Pre-registration drafted (\texttt{files/prereg/C1-credit-estimators.md}); driver not yet built.

\subsection{D1. Breakdown point of verification (edge item 5)}
\textbf{For.} Designing against persuasive-but-wrong ideas.\\
\textbf{What.} Inject a fraction $\varepsilon$ of sabotaged messages (convincing prose, bad code, honest low
evidence) into a small population and find the $\varepsilon$ at which population skill collapses, with and without
execution-based verification before adoption.\\
\textbf{Why.} In robust statistics the breakdown point is the fraction of contamination an estimator tolerates; the
organization is the estimator and bad messages are contamination. Byzantine-robustness papers report thresholds for
consensus rules and memory-poisoning papers report attack rates, but no breakdown curve for execution-based
verification in a learning population was found.\\
\textbf{How.} Population of 8 students in one group (so interference is real and we use the cluster framing), $G=20$
generations, $\varepsilon\in\{0,0.1,0.2,0.3,0.5\}$ assigned by Bernoulli sabotage of messages with known
probability (so exposure probabilities are known, Wager Thm~12.1), verification on/off as the second factor, 10
population seeds per cell for the pilot. Outcome: population mean held-out self-play at $G$, paired on deals.
Breakdown point $\varepsilon^*$: smallest $\varepsilon$ at which the mean falls below $70\%$ of its $\varepsilon=0$
value. Variance by the finite-population estimator with the PSD-projected dependency graph; a permutation test for
the sharp null of no effect of $\varepsilon$ (Wager Thm~11.1). The per-message influence from C1 is the influence
function of the same organization.\\
\textbf{Decision rule.} Claim ``verification raises the breakdown point'' only if $\varepsilon^*_{\text{on}}-
\varepsilon^*_{\text{off}}\ge0.2$ with the bootstrap interval excluding 0, in both model blocks.\\
\textbf{Cost.} $8\times20\times5\times2\times10=16{,}000$ local calls at the pilot size: too many; the pilot uses
4 students, $G=10$, 5 seeds, and the null stub first. Scale after the pilot variance is known.\\
\textbf{Status.} Not yet pre-registered; next after C1.

\subsection{A1. Transmission fidelity and the accumulation prediction (edge item 4)}
\textbf{For.} The formal model that predicts when a population accumulates.\\
\textbf{What.} Send a student a known-good artifact (Piers, $v=17$) in one of three media (code only, prose only,
both); the student makes one revision; measure its self-play. Copy loss $\alpha=v-\text{mean}$; copy variance
$\beta$ from the spread. Henrich's model says mean skill changes per generation by $-\alpha+\beta(0.577+\ln N)$,
$N$ the number of candidates ranked at selection. The sign predicts whether a population improves; we then run a
small population and check the sign. A round trip (code $\to$ document $\to$ fresh reimplementation $\to$
cross-play with the original) asks whether the document is a sufficient statistic of the code for compatibility.\\
\textbf{Why.} Fidelity, not population size, decides accumulation in cultural-evolution theory; no measurement for
LLM agents and no code-versus-prose comparison was found; and the answer tells us what the corpus must store.\\
\textbf{How.} 30 sampling seeds per medium, same deals; report the full distribution per medium, $\hat\alpha,
\hat\beta$ with bootstrap intervals, the predicted sign, and the observed population trajectory over 10 generations
at $N=4$.\\
\textbf{Decision rule.} The model is supported if the predicted sign matches the observed sign in at least 8 of 10
population seeds in both blocks.\\
\textbf{Cost.} 90 calls for the parameters plus $4\times10\times10$ for the check; one night; \$0.\\
\textbf{Status.} Not yet pre-registered.

\subsection{B1. Verification as a sequential test (edge item 5, cost side)}
\textbf{For.} The price of verification and a principled stopping rule.\\
\textbf{What.} Each adoption decision is a paired test on per-deal differences $d_k$ with mean $\delta$ and SD
$\sigma$ (about 3 to 4 for similar bots). Wald's SPRT keeps playing deals until the evidence is decisive; expected
deals $\approx 2\sigma^2\log(1/\alpha)/\delta^2$ at error rates $\alpha,\beta$.\\
\textbf{Why.} Verification is our answer to Jha et al.\ 2026 (social learning net-negative at matched cost) and to
sabotage, and it costs engine time; the SPRT gives the cost as a function of the effect size we care about instead of
a magic $N$. No LLM is involved, so this runs today at \$0.\\
\textbf{How.} Fixed $N\in\{50,100,200,400\}$ versus SPRT at matched error rates on a bank of candidate/incumbent
pairs with known $\delta$; measure false adoptions, missed improvements, deals used.\\
\textbf{Decision rule.} Adopt SPRT as the default if it matches fixed-$N$ error rates with at least $30\%$ fewer
deals at $\delta=1$.\\
\textbf{Status.} Runs on existing anchors; one afternoon.

\subsection{Supporting experiments (cheap, not headline)}
\begin{itemize}
\item \textbf{E1 Winner's curse.} Selecting the best of $n$ noisy evaluations inflates its score; correct with
Efron's Tweedie shrinkage and test whether corrected selection accumulates faster at the same evaluation budget.
\item \textbf{F1 Adaptive evaluation allocation.} Spend verification deals where predicted importance is high with
inverse-probability weights (the Online TASS idea); compare to uniform on deals-to-decision.
\item \textbf{J1 Off-policy evaluation of routing rules.} From one randomized-delivery run, estimate the value of
several routing rules offline with weighted importance sampling and the Wager ch.~15 doubly robust estimator; validate
against two rules actually run. This is the fast-iteration multiplier.
\item \textbf{K1 Designed delivery.} Orthogonal arrays instead of coin flips for identifiable interactions; folded
into C1 Arm~2.
\item \textbf{Infrastructure M, N, O.} Three disjoint deal sets with a generalization gap; null-population
calibration; the innovation base rate $\varepsilon$ (fraction of revisions that improve, with a credible interval)
under three contexts. These are built before any paid or long run.
\end{itemize}

\paragraph{Deferred to post-funding.} Groups competing for compute, migration, selection rules (HGM clade credit,
FunSearch islands), the horizon experiment, mean-field theory of the skill distribution. All are implemented as
frozen, tested policies in the harness and are not extended now.

\section{Novelty and reuse}\label{sec:novelty}

From the arXiv checks of 2026-10-06 (venues not checked separately). Credit: message-level credit is crowded
(Shapley-plus-process-reward, Yang et al.\ 2025; counterfactual replay for failure localization, Li et al.\ 2026;
Aumann--Shapley at scale, Tang et al.\ 2026); we adopt Yang's formalism as the oracle and SHARP's ablation as the
everyday estimator; the randomized-delivery regression checked against exact replay was not found. Breakdown point:
Byzantine thresholds (Liu et al.\ 2026; Lee et al.\ 2026) and memory-poisoning gates (MAPLE-Guard 2026) exist; the
curve for execution-based verification in a learning population was not found. Fidelity: the LLM telephone-game and
cultural-evolution line exists (Perez et al.\ 2024 twice; Vallinder and Hughes 2024; POLIS 2025); Henrich's
parameterization and the code-versus-prose comparison were not found. Defensible phrasings are recorded in
\texttt{files/prereg/novelty-notes.md}.

\section{Build status}\label{sec:status}

Built overnight on 2026-10-05/06 and verified by an independent review: Hanabi engine adapter reproducing published
anchor scores (Piers 16.99, IGGI 15.86, Flawed 0 over 1,000 games, byte-identical to the raw engine); sandboxed
execution of LLM-written bots; self-play and cross-play evaluators; stats (IQM, stratified bootstrap, paired deltas,
Higher Criticism with an empirical null); artifact store with harness-signed evidence and a touch log; LLM cache with
record, replay and strict-replay modes; runner with atomic checkpoints and byte-identical resume; stub backend; 111
tests. Four defects found by review (sandbox escape, soft budget cap, bare \texttt{except} swallowing the timeout,
unpaired HGM labels) and the warm-start race are being fixed, together with the OpenAI-compatible local backend
(sampling seed in the cache key), deal separation, null calibration and the innovation-rate estimator.

Local model: Qwen3.6-35B-A3B via Ollama on the M4 Pro: 71 tokens per second, 30 seconds for a 1,200-token bot,
runnable code on the first try. Under Ollama, identical seeded temperature-0 requests gave different outputs (its runner ignores the
settings); under the MLX server three identical requests were bit-identical, so MLX is the backend for exact replay.

\paragraph{Keep, freeze, add.} Keep the core above; freeze the population layer (groups, migration, allocation,
selection, OpenEvolve adapter) as tested but unused; add one driver for the single-student design with exhaustive
delivery subsets, designed patterns and the pre-registration metrics computed automatically.

\section{Next steps, at agent speed}\label{sec:next}

\begin{enumerate}
\item Commit the staged spec, CLAUDE.md, pre-registrations and this document. Review the fix round when it lands;
run the suite.
\item Resolve local-model determinism (parallelism 1, or MLX). One live revise through the agent step; record
seconds per call and the runnable rate.
\item Enrico marks up C1 (estimand: intent-to-treat or not; thresholds). Revise; draft D1 and A1 in the same
format.
\item Spawn the single-student driver build on a cheap model with the C1 pilot config. Run the C1 pilot (5 seeds) on
the stub, then on the local model, overnight. B1 runs on the anchors the same day at \$0.
\item From the pilot SD, fix $R$; run C1 in full overnight. In parallel, the D1 pilot at reduced size and the null
calibration.
\item C1 figure and table with null bands; ledger row filled in, including if negative. Decide D1 full run. A1.
\item For the preseed deck: two figures (C1, D1), the defensible phrasings, the methods paragraph, and one line on
what funding buys (the factorial over organization rules and population scale).
\item Send Diego the shared evaluation grammar and the split; agree the local model.
\end{enumerate}

\section{Decisions taken on 2026-10-06}
\begin{itemize}
\item C1 estimand: intent-to-treat is primary (delivery is the treatment; verification and adoption are inside the
effect); the Shapley value over delivered subsets is reported alongside; the effect among adopters is reported as
a secondary quantity via delivery as an instrument (Wager ch.~10, one-sided noncompliance), not used in decision
rules.
\item Thresholds as drafted: C1 RMSE $\le2$ points and Spearman IQM $\ge0.9$; D1 collapse at $70\%$ of the
uncontaminated mean and a $0.2$ difference in breakdown points.
\item Order: C1, then D1, then A1; B1 runs on anchors whenever the engine is idle.
\item Replication block: Gemma~4 26B-A4B as the second local model.
\item Local backend: the MLX server (seeded, temperature 0, bit-reproducible, about 65 tokens per second); Ollama's
runner for this model ignores seed and temperature and is not used for replay.
\end{itemize}

\end{document}
